% capitulos/cap5/markdown.tex -- §5.8: pipeline de conteúdo Markdown.
% Parte do Capítulo 5 (capitulos/cap5_arquitetura.tex); os rótulos \label e os números de seção não mudaram.

\section{Pipeline de Conteúdo Markdown}\label{sec:arq-markdown}

A renderização de um documento usa o componente \arq{MarkdownRenderer} (\arq{src/components/markdown/MarkdownRenderer.js}, merge M-009), compartilhado pela página pública e pela pré-visualização do editor, para garantir que o autor veja o que o leitor verá. Como o Markdown vem de autores que não são de confiança (\adr{ADR-023}), a cadeia é restritiva por construção:

\begin{enumerate}
  \item \arq{react-markdown} converte o Markdown em elementos React, \emph{sem HTML bruto} (RN11): sem \arq{rehype-raw}, o HTML escrito no texto aparece como texto;
  \item \arq{remark-gfm} acrescenta tabelas, tachado, listas de tarefas, links automáticos e notas de rodapé (RF11; \decisao{D-08});
  \item \arq{remark-math} reconhece \texttt{\$...\$}, \texttt{\$\$...\$\$} e o bloco de código de linguagem \texttt{math} (RF09);
  \item \arq{rehype-katex} renderiza as fórmulas com \texttt{trust} desligado --- \texttt{\textbackslash href}, \texttt{\textbackslash url}, \texttt{\textbackslash include}\allowbreak\texttt{graphics} e \texttt{\textbackslash html*} não geram marcação ---, \texttt{maxSize} de 30~em, \texttt{maxExpand} de 1000 e \texttt{strict} em \texttt{ignore}; uma fórmula que o KaTeX recusa aparece como texto-fonte, na cor de erro, com o rótulo ``Fórmula inválida:'' e a mensagem original no atributo \texttt{title}, sem derrubar a página;
  \item \arq{rehype-highlight} (sobre o \arq{lowlight} e o \arq{highlight.js}) marca os blocos de código com classes \texttt{hljs-*}, somente quando o bloco declara a linguagem: não há detecção automática, e uma linguagem desconhecida fica sem realce, sem erro (RF10; \decisao{D-07}). Foi escolhido por ser síncrono (o mesmo componente serve ao servidor e à pré-visualização), por emitir só classes (as cores ficam nos tokens) e por não usar WebAssembly nem \texttt{eval}.
\end{enumerate}

\textbf{Segurança.} Vale a transformação padrão de URLs do \arq{react-markdown}, que só admite \texttt{http}, \texttt{https}, \texttt{mailto}, \texttt{irc}, \texttt{ircs}, \texttt{xmpp} e caminhos relativos; \texttt{javascript:}, \texttt{vbscript:}, \texttt{data:} e \texttt{file:} são descartados em links, links de referência, links automáticos e imagens, inclusive com maiúsculas, entidades ou caracteres de controle no esquema. Um link ou uma imagem cuja URL foi descartada \emph{deixa de ser link ou imagem} e vira só o texto. As imagens carregam sob demanda e sem enviar o endereço da página ao servidor da imagem (\texttt{loading} \texttt{lazy} e \texttt{referrerPolicy} \texttt{no-referrer}); uma imagem \texttt{http:} passa pela transformação, mas a CSP a barra (\adr{ADR-022}). O tamanho do corpo (\arq{MAX_BODY_BYTES}, \adr{ADR-023}) é o que limita o número de fórmulas e o tamanho dos blocos de código por documento.

\textbf{Estilo e tema.} A folha do KaTeX é importada uma vez, \emph{pelo próprio} \arq{MarkdownRenderer}, do pacote \arq{katex} (declarado em \texttt{dependencies}), nunca de uma CDN; o Next.js a entrega só às rotas que usam o componente, de modo que o login e o painel não a carregam. Toda cor vem dos tokens de tema (Seção~\ref{sec:arq-tema}): o realce de código usa as variáveis \arq{--vitrine-code-*}, com variante noturna sob a classe \texttt{.dark}, e os links usam \arq{--vitrine-link}, cujo padrão é a cor do texto, sempre sublinhados (Seção~\ref{sec:arq-contraste}).

\textbf{Acessibilidade.} Tabelas e blocos de código rolam dentro do próprio quadro e recebem foco pelo teclado; as caixas das listas de tarefas são só de leitura e dizem o estado por extenso (``Tarefa concluída'' e ``Tarefa pendente''); as fórmulas trazem MathML para leitores de tela; e a seção de notas de rodapé tem os textos em português.

\textbf{Linguagens do realce.} O conjunto é o \texttt{common} do \arq{lowlight}, com 37 linguagens (entre elas C, C++, Java, JavaScript, TypeScript, Python, R, Rust, Go, SQL, Bash, JSON, YAML, HTML e CSS), acrescido de LaTeX, MATLAB, Julia, Haskell e Fortran, que o produto garante por serem comuns no ensino de ciências e de computação (decisão do Arquiteto no P-007, revogável pelo Engenheiro Chefe). A UI as registrou no merge M-011, importando as gramáticas de \texttt{highlight.js/lib/languages/*}: são 42 linguagens, com os apelidos \texttt{tex}, \texttt{hs}, \texttt{f90} e \texttt{f95} que o \arq{highlight.js} já traz, e cada uma tem um teste. O \arq{highlight.js} (BSD-3-Clause, fixado em \texttt{\~{}11.11.2}, para que o projeto tenha uma só cópia, a mesma do \arq{lowlight}) e o \arq{lowlight} (MIT, cujo conjunto \texttt{common} o renderizador estende, em vez de substituí-lo) passaram a ser dependências diretas. COBOL e Scala, por exemplo, seguem sem realce. Outras linguagens entram sob pedido, sem RFC: basta importar a gramática e acrescentá-la a \texttt{HIGHLIGHT\_LANGUAGES}, com um teste.

\textbf{Links internos.} O componente não reescreve links. Um link relativo é resolvido pelo navegador contra a URL da página --- \texttt{outro-documento} aponta para um irmão, e \texttt{../outro} sobe um nível ---, e um link absoluto começa no espaço do autor, em \texttt{/\{handle\}/\ldots} (\adr{ADR-024}). Em consequência, renomear um \en{slug} quebra os links que o citam, e não há redirecionamento nesta versão (Seção~\ref{sec:res-limitacoes}).

\textbf{Desempenho.} Importado por um Server Component, o componente não envia JavaScript ao navegador; dentro de um Componente Cliente (a pré-visualização do editor), acrescenta cerca de 172~kB (\en{gzip}) apenas à página do editor, que deve carregá-lo sob demanda e adiar a renderização em textos longos (\texttt{useDeferredValue}).
